Most prior approaches to cost-constrained BO occur in the \textit{grey-box} setting, in which additional information about the objective is available. 
\textit{Multi-fidelity} BO is a widely studied setting in which fidelity parameters $s \in [0, 1]^m$, such as iteration count or grid size, are assumed to be a low-accuracy approximation of high-fidelity evaluations~\citep{forrester2007multi, kandasamy2017multi, poloczek2017multi, wu2019practical}. 
Increasing $s$ increases the accuracy at the expense of run time. 
In addition, Multi-fidelity methods are often application-specific. 
For example, Hyperband~\citep{li2017hyperband, falkner2018bohb, klein2016fast, klein2017fastbo} cheaply train many neural network configurations for a few epochs, and then prunes unpromising configurations. 

In \textit{multi-task} BO ~\citep{swersky2013multi}, hyperparameter optimization is first run on cheaper instances before considering more expensive ones. 
\citet{swersky2013multi} introduce a cost-constrained, multi-task variant of entropy search to speed-up optimization of logistic regression and latent Dirichlet allocation. 
Cost information is input as a set of cost preferences (e.g., a parameter $\bx_1$ is more expensive than a parameter $\bx_2$) by \citet{abdolshah2019costaware}, who develop a multi-objective, constrained BO method that evaluates cheap points before expensive ones, as determined by the cost preferences, to find feasible, low-cost solutions. 
These methods outperform their black-box counterparts by evaluating cheap proxies or cheap points before selecting expensive evaluations, which is accomplished by leveraging additional cost information inside the optimization routine. 

While all these methods demonstrate strong performance, they sacrifice generality and do not apply to black-box BO. 
Moreover, by relying on parallel resources, these techniques target time efficiency rather than compute time, cost, or energy efficiency. We also note that cost-efficiency has been modestly studied in the setting of active search, in which a active search is run with the constraint that the total cost of all queries must be less than $\tau$ \cite{jiang2019cost}.

Recent developments in nonmyopic BO account for the impact of future evaluations and are thus able to make better decisions~\citep{lam2016bayesian, lam2017lookahead, frazier2018, yue2019lookahead, lee2020efficient, jiang2020efficient}. These methods typically frame Bayesian optimization as a Markov decision process whose horizon is precisely the iteration budget. It is along these lines that we tackle the problem of nonmyopic, cost-constrained BO.   